\documentclass[11pt,a4paper]{article}

% ---------- Packages ----------
\usepackage[utf8]{inputenc}
\usepackage[T1]{fontenc}
\usepackage[english]{babel}
\usepackage{lmodern}
\usepackage{geometry}
\geometry{margin=2.5cm}
\usepackage{graphicx}
\usepackage{longtable}
\usepackage{array}
\usepackage{booktabs}
\usepackage{xcolor}
\usepackage{listings}
\usepackage{enumitem}
\usepackage{hyperref}
\usepackage{titlesec}
\usepackage{fancyhdr}

% ---------- Colors ----------
\definecolor{codebg}{RGB}{245,245,245}
\definecolor{codeframe}{RGB}{210,210,210}
\definecolor{linkblue}{RGB}{0,90,160}

\hypersetup{
    colorlinks=true,
    linkcolor=linkblue,
    urlcolor=linkblue,
    citecolor=linkblue,
    pdftitle={NoFluxo UNB - Project Documentation},
    pdfauthor={Squad 03 - MDS 2025/1, FGA/UnB}
}

% ---------- Code listing style ----------
\lstset{
    basicstyle=\ttfamily\small,
    backgroundcolor=\color{codebg},
    frame=single,
    rulecolor=\color{codeframe},
    breaklines=true,
    columns=fullflexible,
    keepspaces=true,
    showstringspaces=false,
    xleftmargin=0.5em,
    xrightmargin=0.5em,
    aboveskip=1em,
    belowskip=1em
}

% ---------- Headers / Footers ----------
\pagestyle{fancy}
\fancyhf{}
\lhead{NoFluxo UNB}
\rhead{Project Documentation}
\cfoot{\thepage}
\renewcommand{\headrulewidth}{0.4pt}

% Convenience macros
\newcommand{\code}[1]{\texttt{#1}}

% Column types for tables
\newcolumntype{L}[1]{>{\raggedright\arraybackslash}p{#1}}

\titlespacing*{\section}{0pt}{1.5em}{0.6em}
\titlespacing*{\subsection}{0pt}{1.1em}{0.4em}

\begin{document}

% ---------- Title ----------
\begin{titlepage}
    \centering
    \vspace*{4cm}
    {\Huge\bfseries NoFluxo UNB\par}
    \vspace{0.6cm}
    {\LARGE Project Documentation\par}
    \vspace{2cm}
    {\large Squad 03 — MDS 2025/1\par}
    {\large (Métodos de Desenvolvimento de Software)\par}
    {\large FGA/UnB\par}
    \vspace{1.5cm}
    {\large Live URL: \url{https://no-fluxo.com}\par}
    \vfill
    {\large \today\par}
\end{titlepage}

\tableofcontents
\newpage

% =====================================================================
\section{Overview}
% =====================================================================

\textbf{NoFluxo} is a full-stack web application that helps students at the
Universidade de Brasília (UnB) visualize and plan their academic curriculum.
Students can upload their academic transcript (PDF), and the system parses it,
matches completed disciplines to the official curriculum, and displays an
interactive flowchart showing their progress.

\medskip
\noindent\textbf{Live URL:} \url{https://no-fluxo.com}

\medskip
\noindent\textbf{Core capabilities:}
\begin{itemize}[noitemsep]
    \item Interactive curriculum flowchart with prerequisite visualization
    \item PDF transcript parsing (browser-side + OCR fallback)
    \item Discipline matching algorithm (transcript vs.\ official curriculum)
    \item Equivalency detection across curriculum versions
    \item AI-powered discipline recommendations (via RAGFlow)
    \item Google OAuth, email/password, and anonymous authentication
    \item Progress tracking with mandatory/elective hour calculations
\end{itemize}

\medskip
\noindent\textbf{Tech Stack:}

\begin{center}
\begin{tabular}{@{}ll@{}}
\toprule
\textbf{Layer} & \textbf{Technology} \\
\midrule
Frontend       & SvelteKit 2 + Svelte 5, Tailwind CSS 4, Vite 7 \\
Backend        & Node.js + TypeScript + Express.js \\
PDF Processing & Python (PyMuPDF + Tesseract OCR) \\
Database       & Supabase (PostgreSQL) \\
AI             & RAGFlow agent \\
Deployment     & Kubernetes (K3s), Docker, GitHub Actions \\
CI/CD          & GitHub Actions \\
\bottomrule
\end{tabular}
\end{center}

\medskip
\noindent\textbf{Team:} Squad 03 — MDS 2025/1 (Métodos de Desenvolvimento de
Software), FGA/UnB.

% =====================================================================
\section{Project Structure}
% =====================================================================

\begin{lstlisting}
2025-1-NoFluxoUNB/
|-- no_fluxo_backend/              # TypeScript/Express backend + Python PDF parsing
|-- no_fluxo_frontend_svelte/      # SvelteKit frontend application
|-- coleta_dados/                  # Data collection/scraping scripts (Python)
|-- documentacao/                  # Project documentation (Portuguese)
|-- kubernetes_docs/               # Kubernetes deployment docs & templates
|-- plans/                         # Planning and strategy documents
|-- test_historicos/               # Test PDF transcripts for parsing
|-- tests-python/                  # Python test suite
|-- scripts/
|   `-- deploy/                    # Deploy scripts for Kubernetes (deploy_local.py)
|-- .github/                       # GitHub Actions CI/CD workflows
|-- k8s.backend.Dockerfile         # Kubernetes backend image
|-- k8s.frontend-svelte.Dockerfile # Kubernetes frontend image
|-- requirements.txt               # Root Python dependencies
|-- README.md                      # Main project README
|-- CONTRIBUTING.md                # Contribution guidelines
|-- COMMIT_GUIDELINES.md           # Git commit conventions
|-- CODE_OF_CONDUCT.md             # Code of conduct
|-- SECURITY.md                    # Security policy
`-- LICENSE                        # GNU GPLv3 License
\end{lstlisting}

% =====================================================================
\section{Backend}
% =====================================================================

\noindent\textbf{Location:} \code{no\_fluxo\_backend/}\\
\textbf{Port:} 3325\\
\textbf{Language:} TypeScript (Express.js) + Python (PDF parsing)

\subsection{Directory Structure}

\begin{lstlisting}
no_fluxo_backend/
|-- src/
|   |-- index.ts                    # Express server entry point
|   |-- interfaces.ts               # TypeScript interfaces
|   |-- supabase_wrapper.ts         # Supabase client wrapper
|   |-- logger.ts                   # Winston logger config
|   |-- utils.ts                    # General utilities
|   |-- controllers/
|   |   |-- fluxograma_controller.ts   # Core flowchart/matching logic
|   |   |-- users_controller.ts        # User registration
|   |   |-- cursos_controller.ts       # Course listing
|   |   |-- materias_controller.ts     # Discipline lookup
|   |   |-- assistente_controller.ts   # AI assistant proxy
|   |   `-- testes_controller.ts       # Debug/test endpoints
|   |-- services/
|   |   |-- ragflow.service.ts         # RAGFlow API client
|   |   `-- ragflow.types.ts           # RAGFlow type definitions
|   `-- utils/
|       |-- controller_logger.ts       # Scoped controller logging
|       |-- text.utils.ts              # Text processing (accents, HTML)
|       `-- ranking.formatter.ts       # AI response -> Markdown formatter
|-- parse-pdf/
|   |-- pdf_parser_final.py            # Main PDF parser (PyMuPDF + Tesseract)
|   |-- pdf_parser_ocr.py              # OCR-only parser
|   |-- requirements.txt               # Python dependencies
|   `-- tests/                         # Python tests
|-- scripts/
|   |-- export_schema.ts               # Database schema export tool
|   |-- setup-tesseract.js             # Tesseract OCR setup
|   `-- migrations/                    # SQL migration files
|-- docs/
|   |-- database_schema.json           # Exported full schema
|   |-- database_tables.md             # Table documentation
|   |-- database_functions.sql         # Stored functions
|   |-- rls_policies.sql               # Row-level security policies
|   `-- supabase_migrations/           # Migration SQL files
|-- package.json
|-- tsconfig.json
`-- Dockerfile
\end{lstlisting}

\subsection{Key Dependencies}
\begin{itemize}[noitemsep]
    \item \textbf{express} — HTTP framework
    \item \textbf{@supabase/supabase-js} — Supabase client
    \item \textbf{axios} — HTTP client (RAGFlow calls)
    \item \textbf{helmet} — Security headers
    \item \textbf{cors} — Cross-origin middleware
    \item \textbf{express-fileupload} — File upload handling
    \item \textbf{express-ws} — WebSocket support
    \item \textbf{winston / morgan} — Logging
\end{itemize}

\subsection{Controllers}

\begin{center}
\begin{longtable}{@{}L{4.2cm}L{2.8cm}L{6cm}@{}}
\toprule
\textbf{Controller} & \textbf{Base Route} & \textbf{Purpose} \\
\midrule
\endhead
\code{fluxograma\_controller.ts}  & \code{/fluxograma}  & Flowchart retrieval, discipline matching, save/delete \\
\code{users\_controller.ts}       & \code{/users}       & User registration (Google OAuth, email) \\
\code{cursos\_controller.ts}      & \code{/cursos}      & Course catalog listing \\
\code{materias\_controller.ts}    & \code{/materias}    & Discipline lookup by code/name \\
\code{assistente\_controller.ts}  & \code{/assistente}  & AI assistant proxy to RAGFlow \\
\code{testes\_controller.ts}      & \code{/testes}      & Debug/testing endpoints \\
\bottomrule
\end{longtable}
\end{center}

\subsection{Core Business Logic: Discipline Matching}

The most important endpoint is \code{POST /fluxograma/casar\_disciplinas}. It runs
a \textbf{4-stage matching algorithm}:

\begin{enumerate}[noitemsep]
    \item \textbf{Primary Matrix Matching} — Matches extracted disciplines against
    the student's curriculum version by code and name.
    \item \textbf{Alternative Matrix Matching} — Checks other curriculum versions of
    the same course.
    \item \textbf{Equivalency Matching} — Searches the equivalency table for matching
    logical expressions.
    \item \textbf{Classification} — Categorizes each discipline by status
    (\code{APR}, \code{CUMP}, \code{MATR}, \code{REP}) and type
    (\code{obrigatoria}, \code{optativa}).
\end{enumerate}

The algorithm handles duplicates (same discipline with different statuses), uses
priority-based status selection, and computes completion metrics (hours,
percentages, IRA).

\subsection{PDF Parsing (Python)}

Located in \code{parse-pdf/}, the Python service extracts structured data from UnB
academic transcript PDFs:

\begin{itemize}[noitemsep]
    \item Uses \textbf{PyMuPDF (fitz)} for text extraction
    \item Falls back to \textbf{Tesseract OCR} for scanned documents
    \item Extracts: course name, curriculum version, disciplines (code, name,
    status, grade, hours), IRA, weighted average, frequency
    \item Handles SIGAA (UnB's academic system) format
\end{itemize}

\subsection{Logging}

Uses Winston with three transports:
\begin{itemize}[noitemsep]
    \item Console (colored output)
    \item \code{logs/error.log} (errors only)
    \item \code{logs/all.log} (all levels)
\end{itemize}

Controller-level scoped logging via \code{ControllerLogger} class:
\code{[Controller][Endpoint] message}.

% =====================================================================
\section{Frontend}
% =====================================================================

\noindent\textbf{Location:} \code{no\_fluxo\_frontend\_svelte/}\\
\textbf{Framework:} SvelteKit 2 + Svelte 5 (runes)\\
\textbf{Styling:} Tailwind CSS 4\\
\textbf{Build Tool:} Vite 7

\subsection{Directory Structure}

\begin{lstlisting}
no_fluxo_frontend_svelte/src/
|-- routes/                          # SvelteKit pages
|   |-- +layout.svelte               # Root layout (auth + navigation)
|   |-- +page.svelte                 # Landing page
|   |-- login/                       # Login page
|   |-- signup/                      # Registration page
|   |-- login-anonimo/               # Anonymous login
|   |-- password-recovery/           # Password recovery
|   |-- auth/callback/               # OAuth callback
|   |-- auth/reset-password/         # Password reset
|   |-- fluxogramas/                 # Course listing page
|   |-- meu-fluxograma/              # User's flowchart
|   |-- meu-fluxograma/[courseName]/ # Dynamic course flowchart
|   |-- upload-historico/            # PDF transcript upload
|   |-- assistente/                  # AI assistant (in development)
|   |-- api/casar-disciplinas/       # SvelteKit API route
|   |-- health/                      # Health check
|   `-- sitemap.xml/                 # SEO sitemap
|
|-- lib/
|   |-- components/
|   |   |-- auth/                   # Login, Signup, Logout, Recovery forms
|   |   |-- fluxograma/             # Flowchart display components
|   |   |-- upload/                 # PDF upload flow components
|   |   |-- layout/                 # Navbar, Sidebar, Footer, Logo, etc.
|   |   |-- home/                   # Landing page sections
|   |   |-- effects/                # AnimatedBackground, GraffitiBackground
|   |   |-- forms/                  # Select, TextInput
|   |   |-- ui/                     # bits-ui wrappers (Button, Card, Dialog, etc.)
|   |   |-- icons/                  # GoogleIcon
|   |   `-- seo/                    # PageMeta, JsonLd
|   |
|   |-- services/
|   |   |-- auth.service.ts         # Supabase authentication
|   |   |-- fluxograma.service.ts   # Flowchart data loading
|   |   |-- upload.service.ts       # PDF upload & processing
|   |   |-- assistente.service.ts   # AI assistant API
|   |   |-- integralizacao.service.ts # Completion calculation
|   |   |-- user.service.ts         # User profile operations
|   |   |-- supabase-data.service.ts # Direct Supabase queries (RLS)
|   |   `-- pdf/
|   |       |-- pdfParser.ts        # PDF parsing orchestrator
|   |       |-- pdfExtractor.ts     # Text extraction (PDF.js)
|   |       |-- pdfDataExtractor.ts # Structured data parsing
|   |       `-- pdfPositionExtractor.ts # Position-based parsing
|   |
|   |-- stores/                     # State management
|   |   |-- auth.ts                 # Auth state (Svelte 4 writable)
|   |   |-- fluxograma.store.svelte.ts # Flowchart state (Svelte 5 runes)
|   |   |-- chatStore.ts            # AI chat state
|   |   |-- uploadStore.ts          # Upload progress state
|   |   |-- modal.ts                # Modal visibility
|   |   |-- navigation.ts           # Navigation state
|   |   `-- theme.ts                # Dark mode
|   |
|   |-- types/                      # TypeScript type definitions
|   |-- schemas/                    # Zod validation schemas
|   |-- guards/                     # Route protection (authGuard.ts)
|   |-- hooks/                      # useBreakpoint, useMediaQuery, useSession
|   |-- supabase/                   # Browser Supabase client
|   |-- server/                     # Server Supabase client
|   |-- utils/                      # Helpers (API, breakpoints, expressions, toast)
|   |-- actions/                    # Svelte actions (clickOutside, inview)
|   |-- factories/                  # Data transformation factories
|   `-- config.ts                   # Environment config
|
|-- app.css                         # Global Tailwind styles
|-- app.d.ts                        # Global type declarations
`-- hooks.server.ts                 # Server hooks (auth, redirects)
\end{lstlisting}

\subsection{Pages \& Routes}

\begin{center}
\begin{longtable}{@{}L{4.5cm}L{2.2cm}L{6cm}@{}}
\toprule
\textbf{Route} & \textbf{Auth Required} & \textbf{Purpose} \\
\midrule
\endhead
\code{/}                         & No  & Landing page (hero, features, about) \\
\code{/login}                    & No  & Email/password login \\
\code{/signup}                   & No  & User registration \\
\code{/password-recovery}        & No  & Password reset request \\
\code{/login-anonimo}            & No  & Anonymous browsing mode \\
\code{/auth/callback}            & No  & Google OAuth callback handler \\
\code{/auth/reset-password}      & No  & Set new password form \\
\code{/fluxogramas}              & Yes & List all available courses/flowcharts \\
\code{/meu-fluxograma}           & Yes & User's personal flowchart overview \\
\code{/meu-fluxograma/[course]}  & Yes & Course-specific interactive flowchart \\
\code{/upload-historico}         & Yes & Upload PDF academic transcript \\
\code{/assistente}               & Yes & AI assistant (in development) \\
\bottomrule
\end{longtable}
\end{center}

\subsection{Key Component Groups}

\noindent\textbf{Flowchart Components} — The main feature of the app:
\begin{itemize}[noitemsep]
    \item \code{FluxogramContainer.svelte} — Interactive grid with pan/zoom (mouse + touch pinch)
    \item \code{SemesterColumn.svelte} — Renders a single semester's subjects in a column
    \item \code{SubjectCard.svelte} — Individual subject card (color-coded by status)
    \item \code{PrerequisiteConnections.svelte} — SVG lines connecting prerequisites
    \item \code{SubjectDetailsModal.svelte} — Detailed subject info on click
    \item \code{PrerequisiteChainDialog.svelte} — Full prerequisite chain visualization
    \item \code{OptativasModal.svelte} — Elective subject selection
    \item \code{ProgressSummarySection.svelte} — Completion summary (mandatory/elective hours)
    \item \code{IntegralizacaoSection.svelte} — Credit requirements progress
\end{itemize}

\noindent\textbf{Upload Components} — PDF processing flow:
\begin{itemize}[noitemsep]
    \item \code{FileDropzone.svelte} — Drag-and-drop PDF upload area
    \item \code{UploadProgress.svelte} — Progress indicator
    \item \code{CourseSelectionModal.svelte} — Course disambiguation
    \item \code{DisciplinaList.svelte} — Matched discipline listing
    \item \code{UploadSuccess.svelte} — Success screen
    \item \code{ProcessingResults.svelte} — Parsing results
\end{itemize}

\noindent\textbf{Layout Components:}
\begin{itemize}[noitemsep]
    \item \code{Navbar.svelte}, \code{Sidebar.svelte}, \code{Footer.svelte}
    \item \code{AppLogo.svelte}, \code{Breadcrumbs.svelte}
    \item \code{LoadingBar.svelte}, \code{SplashScreen.svelte}
    \item \code{GlassContainer.svelte} — Glassmorphism container
    \item \code{GradientCTAButton.svelte} — Animated CTA button
\end{itemize}

\noindent\textbf{Landing Page Sections:}
\begin{itemize}[noitemsep]
    \item \code{HeroSection.svelte} — Main banner
    \item \code{ComoFuncionaSection.svelte} — ``How it works'' steps
    \item \code{ProntoParaOrganizarSection.svelte} — Call to action
    \item \code{SobreNosSection.svelte} — About the team
\end{itemize}

\subsection{State Management}

The app uses a mix of Svelte 4 writable stores and Svelte 5 runes:

\begin{center}
\begin{longtable}{@{}L{4.8cm}L{3cm}L{5cm}@{}}
\toprule
\textbf{Store} & \textbf{Type} & \textbf{Purpose} \\
\midrule
\endhead
\code{auth.ts}                      & Svelte writable & Auth state, user data, localStorage sync \\
\code{fluxograma.store.svelte.ts}   & Svelte 5 runes  & Course data, zoom, selected subjects \\
\code{uploadStore.ts}               & Svelte writable & Upload progress, matched disciplines \\
\code{chatStore.ts}                 & Svelte writable & AI chat messages \\
\code{modal.ts}                     & Svelte writable & Modal open/close state \\
\code{navigation.ts}                & Svelte writable & Navigation state \\
\code{theme.ts}                     & Svelte writable & Dark/light mode \\
\bottomrule
\end{longtable}
\end{center}

\subsection{Services Layer}

\begin{center}
\begin{longtable}{@{}L{4.8cm}L{8cm}@{}}
\toprule
\textbf{Service} & \textbf{Purpose} \\
\midrule
\endhead
\code{auth.service.ts}            & Supabase Auth (login, signup, OAuth, anonymous, password) \\
\code{fluxograma.service.ts}      & Load course data and build flowchart models \\
\code{upload.service.ts}          & Browser-side PDF parsing + discipline matching \\
\code{supabase-data.service.ts}   & Direct Supabase queries with RLS (courses, matrices, subjects) \\
\code{integralizacao.service.ts}  & Calculate mandatory/elective hour completion \\
\code{assistente.service.ts}      & Proxy to AI assistant backend \\
\code{pdf/pdfParser.ts}           & PDF.js-based browser-side transcript parsing \\
\bottomrule
\end{longtable}
\end{center}

\subsection{Key Data Flows}

\noindent\textbf{PDF Upload Flow:}
\begin{lstlisting}
FileDropzone -> uploadService.parsePdfLocally() -> PDF.js parsing (browser)
-> Extract disciplines -> POST /api/casar-disciplinas -> Backend matching
-> Display results -> User confirms -> saveFluxogramaData() -> Navigate to flowchart
\end{lstlisting}

\noindent\textbf{Flowchart View Flow:}
\begin{lstlisting}
Route loads -> authStore user data -> fluxogramaService.getCourseData()
-> Query Supabase (RLS) -> Build CursoModel -> fluxogramaStore
-> Render FluxogramContainer -> Semester columns + subject cards + SVG connections
\end{lstlisting}

\subsection{Path Aliases}

\begin{lstlisting}
$components -> src/lib/components
$lib        -> src/lib
$stores     -> src/lib/stores
$types      -> src/lib/types
$services   -> src/lib/services
$utils      -> src/lib/utils
\end{lstlisting}

% =====================================================================
\section{Database}
% =====================================================================

\noindent\textbf{Platform:} Supabase (PostgreSQL)\\
\textbf{URL:} \code{https://lijmhbstgdinsukovyfl.supabase.co}

\subsection{Tables}

\begin{center}
\begin{longtable}{@{}L{3.5cm}rrL{5.5cm}@{}}
\toprule
\textbf{Table} & \textbf{Rows} & \textbf{Size} & \textbf{Purpose} \\
\midrule
\endhead
\code{users}              & 162    & 136 KB  & User accounts (linked to Supabase auth) \\
\code{dados\_users}        & 99     & 600 KB  & User flowchart data and current semester \\
\code{cursos}             & 102    & 64 KB   & UnB courses (Administração, Engenharia, etc.) \\
\code{matrizes}           & 235    & 96 KB   & Curriculum versions per course \\
\code{materias}           & 25,624 & 8.4 MB  & Discipline catalog (code, name, hours, syllabus) \\
\code{materias\_por\_curso} & 65,955 & 13.1 MB & Discipline-to-curriculum mapping (with level) \\
\code{pre\_requisitos}     & 14,573 & 2.3 MB  & Prerequisite relationships \\
\code{co\_requisitos}      & 185    & 80 KB   & Corequisite relationships \\
\code{equivalencias}      & 19,445 & 4 MB    & Subject equivalencies across curricula \\
\bottomrule
\end{longtable}
\end{center}

\subsection{Key Relationships}

\begin{lstlisting}
auth.users (Supabase Auth)
    `-- users.auth_id
         `-- dados_users.id_user     (user's saved flowchart)

cursos
    `-- matrizes.id_curso            (curriculum versions)
         `-- materias_por_curso.id_matriz
              `-- materias.id_materia (disciplines in curriculum)

materias
    |-- pre_requisitos (self-referential: materia -> materia_requisito)
    |-- co_requisitos  (self-referential: materia -> materia_corequisito)
    `-- equivalencias  (materia + curso + matriz)
\end{lstlisting}

\subsection{Key Fields}
\begin{itemize}[noitemsep]
    \item \code{materias\_por\_curso.nivel} — \code{0} = elective, \code{1-8} = mandatory at that semester
    \item \code{equivalencias.expressao\_logica} — JSONB with logical expressions for equivalency rules
    \item \code{matrizes.curriculo\_completo} — Unique curriculum identifier string
    \item \code{matrizes.ch\_obrigatoria\_exigida} / \code{ch\_optativa\_exigida} — Required hours for graduation
\end{itemize}

\subsection{Views}
\begin{itemize}[noitemsep]
    \item \textbf{\code{vw\_creditos\_por\_matriz}} — Joins matrizes $\to$ cursos $\to$
    materias\_por\_curso $\to$ materias to compute credit totals per curriculum version.
\end{itemize}

\subsection{Custom Functions}
\begin{itemize}[noitemsep]
    \item \code{export\_schema()} — Exports complete database schema as JSONB
    \item \code{calcular\_creditos\_por\_curso(id\_curso)} — Calculates credits (1 credit = 15 hours)
    \item \code{atualizar\_creditos\_cursos()} — Batch update credits for all courses
\end{itemize}

\subsection{Row Level Security (RLS)}

All tables have RLS enabled. Policies:
\begin{itemize}[noitemsep]
    \item \textbf{Academic data} (cursos, materias, matrizes, pre\_requisitos,
    co\_requisitos, equivalencias): \textbf{Public read} for authenticated users.
    \item \textbf{users}: Users can only read/update their \textbf{own} record
    (\code{WHERE auth\_id = auth.uid()}).
    \item \textbf{dados\_users}: Full CRUD restricted to the \textbf{owning user} only.
\end{itemize}

% =====================================================================
\section{Data Collection}
% =====================================================================

\noindent\textbf{Location:} \code{coleta\_dados/}

\subsection{Scraping Scripts}

Python scripts that scrape data from UnB's SIGAA (academic information system):

\begin{center}
\begin{longtable}{@{}L{5.5cm}L{7cm}@{}}
\toprule
\textbf{Script} & \textbf{Purpose} \\
\midrule
\endhead
\code{scraping\_cursos.py}        & Scrape course information \\
\code{scraping\_atual.py}         & Current course data \\
\code{scraping\_ementa.py}        & Course syllabi \\
\code{scraping\_equivalencias.py} & Subject equivalencies \\
\code{scraping\_sigaa.py}         & SIGAA system integration \\
\code{extrair\_turmas\_sigaa.py}   & Extract class sections \\
\code{integracao\_banco.py} (51 KB) & Main database integration layer \\
\code{insere\_departamento.py}    & Insert department data \\
\code{insere\_tipocurso.py}       & Insert course type data \\
\code{confere\_estruturas.py}     & Verify data structures \\
\code{formatar\_para\_ragflow.py}  & Format data for RAGFlow AI \\
\bottomrule
\end{longtable}
\end{center}

\subsection{Data Files}
\begin{itemize}[noitemsep]
    \item \code{dados/cursos-de-graduacao.json} — 102+ UnB courses with metadata
    \item \code{dados/materias/turmas\_depto\_*.json} — Class sections organized by department (20+ files)
\end{itemize}

\subsection{Pipeline}

\begin{lstlisting}
SIGAA website -> Scraping scripts -> JSON data files -> integracao_banco.py -> Supabase DB
                                                      -> formatar_para_ragflow.py -> RAGFlow AI
\end{lstlisting}

% =====================================================================
\section{Infrastructure \& Deployment}
% =====================================================================

\subsection{Kubernetes (Production)}

Production runs on a K3s cluster. The deploy workflow is:

\begin{lstlisting}
Developer runs scripts/deploy/deploy_local.py backend|frontend
  -> Docker build (multi-stage, local or buildx)
  -> Push to registry.kubernetes.crianex.com
  -> Deploy API resolves image digest
  -> Deploy API creates K8s manifests and applies via kubectl
  -> Service available at crianex.com domain
\end{lstlisting}

\noindent\textbf{Cluster details:}
\begin{itemize}[noitemsep]
    \item \textbf{Platform:} K3s cluster at \code{kubernetes.crianex.com}
    \item \textbf{Ingress:} Traefik (handles SSL/TLS via cert-manager + Let's Encrypt)
    \item \textbf{Registry:} Private Docker registry (\code{registry.kubernetes.crianex.com})
    \item \textbf{Deploy API:} REST API at \code{deploy.kubernetes.crianex.com}
    \item \textbf{Monitoring:} Prometheus + Grafana + Loki
\end{itemize}

\noindent\textbf{Production domains:}
\begin{itemize}[noitemsep]
    \item Frontend: \code{https://no-fluxo.crianex.com}
    \item Backend API: \code{https://api-nofluxo.crianex.com}
\end{itemize}

\subsection{Deploy Scripts}

Located in \code{scripts/deploy/}:

\begin{center}
\begin{longtable}{@{}L{3.5cm}L{9cm}@{}}
\toprule
\textbf{File} & \textbf{Purpose} \\
\midrule
\endhead
\code{deploy\_local.py}  & Main deploy script. Builds Docker images, pushes to registry, deploys via API \\
\code{deploy\_config.py} & App configuration (names, ports, domains, env vars) for backend and frontend \\
\code{requirements.txt} & Python dependencies (\code{requests}, \code{python-dotenv}) \\
\code{.deploy/}         & Cached build IDs to avoid redundant rebuilds \\
\bottomrule
\end{longtable}
\end{center}

\subsection{Dockerfiles}

\begin{center}
\begin{longtable}{@{}L{3.8cm}L{1.8cm}L{2.5cm}L{4.5cm}@{}}
\toprule
\textbf{File} & \textbf{Target} & \textbf{Base} & \textbf{Description} \\
\midrule
\endhead
\code{k8s.backend.Dockerfile}         & Backend  & Node 20 Alpine   & Multi-stage build, non-root user, health check, runs \code{node dist/index.js} \\
\code{k8s.frontend-svelte.Dockerfile} & Frontend & Node 20 Bullseye & Multi-stage build, pnpm, adapter-node for SSR, build-time env vars \\
\bottomrule
\end{longtable}
\end{center}

\subsection{Reference Documentation}

The \code{kubernetes\_docs/} folder contains reference documentation and templates
for the Kubernetes cluster setup, including registry auth guides, monitoring
dashboards, and generic Dockerfile templates.

\subsection{CORS Configuration}

Allowed origins:
\begin{itemize}[noitemsep]
    \item \code{https://www.no-fluxo.com}
    \item \code{https://no-fluxo.com}
    \item \code{http://localhost:3000}
    \item \code{http://localhost:5000}
\end{itemize}

% =====================================================================
\section{CI/CD \& Testing}
% =====================================================================

\subsection{GitHub Actions Workflows}

\begin{center}
\begin{longtable}{@{}L{4.5cm}L{8cm}@{}}
\toprule
\textbf{Workflow} & \textbf{Purpose} \\
\midrule
\endhead
\code{all-tests.yml}            & Run TypeScript + Python tests \\
\code{typescript-tests.yml}     & ESLint + type-check + Jest \\
\code{python-tests.yml}         & pytest on Python 3.9, 3.10, 3.11 \\
\code{security-and-quality.yml} & Security scans + code quality checks \\
\code{pipelineCI.yml}           & CI pipeline \\
\bottomrule
\end{longtable}
\end{center}

\subsection{Testing Tools}

\begin{center}
\begin{longtable}{@{}L{3cm}L{3cm}L{5cm}@{}}
\toprule
\textbf{Tool} & \textbf{Language} & \textbf{Purpose} \\
\midrule
\endhead
Jest       & TypeScript & Backend unit tests \\
Vitest     & TypeScript & Frontend unit tests \\
Playwright & TypeScript & E2E integration \\
pytest     & Python     & PDF parser tests \\
\bottomrule
\end{longtable}
\end{center}

\subsection{Quality Tools}
\begin{itemize}[noitemsep]
    \item \textbf{TypeScript:} ESLint, Prettier, TypeScript strict mode
    \item \textbf{Python:} flake8, black, isort, mypy, bandit, safety
    \item \textbf{Coverage:} Codecov integration
\end{itemize}

% =====================================================================
\section{API Reference}
% =====================================================================

\subsection*{Backend API (Express, port 3325)}

\subsubsection*{Health}
\begin{center}
\begin{longtable}{@{}L{2cm}L{3cm}L{6cm}@{}}
\toprule
\textbf{Method} & \textbf{Endpoint} & \textbf{Description} \\
\midrule
\endhead
GET & \code{/}       & Server info + endpoint listing \\
GET & \code{/health} & Kubernetes health probe \\
\bottomrule
\end{longtable}
\end{center}

\subsubsection*{Users (\code{/users})}
\begin{center}
\begin{longtable}{@{}L{1.5cm}L{4.5cm}L{3.5cm}L{3cm}@{}}
\toprule
\textbf{Method} & \textbf{Endpoint} & \textbf{Body} & \textbf{Description} \\
\midrule
\endhead
POST & \code{/users/register-user-with-google}  & \code{\{ email, nome\_completo \}} & Register via Google OAuth \\
GET  & \code{/users/get-user-by-email}           & Query: \code{email}               & Get user profile \\
POST & \code{/users/registrar-user-with-email}   & \code{\{ email, nome\_completo \}} & Register with email \\
\bottomrule
\end{longtable}
\end{center}

\subsubsection*{Courses (\code{/cursos})}
\begin{center}
\begin{longtable}{@{}L{1.5cm}L{4cm}L{6cm}@{}}
\toprule
\textbf{Method} & \textbf{Endpoint} & \textbf{Description} \\
\midrule
\endhead
GET & \code{/cursos/all-cursos} & List all courses with credit info \\
\bottomrule
\end{longtable}
\end{center}

\subsubsection*{Disciplines (\code{/materias})}
\begin{center}
\begin{longtable}{@{}L{1.5cm}L{4.5cm}L{4cm}L{3cm}@{}}
\toprule
\textbf{Method} & \textbf{Endpoint} & \textbf{Body} & \textbf{Description} \\
\midrule
\endhead
GET  & \code{/materias/materias-name-by-code} & \code{\{ codes: string[] \}}             & Get names by codes \\
POST & \code{/materias/materias-from-codigos} & \code{\{ codigos: string[], id\_curso \}} & Get detailed subject info \\
\bottomrule
\end{longtable}
\end{center}

\subsubsection*{Flowchart (\code{/fluxograma})}
\begin{center}
\begin{longtable}{@{}L{1.5cm}L{4.5cm}L{3.5cm}L{3.5cm}@{}}
\toprule
\textbf{Method} & \textbf{Endpoint} & \textbf{Body / Query} & \textbf{Description} \\
\midrule
\endhead
GET    & \code{/fluxograma/fluxograma}             & Query: \code{nome\_curso}              & Get full curriculum flowchart \\
POST   & \code{/fluxograma/casar\_disciplinas}      & Extracted PDF data                    & Match transcript to curriculum \\
POST   & \code{/fluxograma/upload-dados-fluxograma} & \code{\{ fluxograma, periodo\_letivo \}} & Save flowchart to user profile \\
DELETE & \code{/fluxograma/delete-fluxograma}       & Headers: Authorization, User-ID       & Remove user's saved flowchart \\
\bottomrule
\end{longtable}
\end{center}

\subsubsection*{AI Assistant (\code{/assistente})}
\begin{center}
\begin{longtable}{@{}L{1.5cm}L{3.5cm}L{3.5cm}L{4cm}@{}}
\toprule
\textbf{Method} & \textbf{Endpoint} & \textbf{Body} & \textbf{Description} \\
\midrule
\endhead
POST & \code{/assistente/analyze} & \code{\{ materia: string \}} & Get AI discipline recommendations \\
GET  & \code{/assistente/health}  & ---                          & Check AI service status \\
\bottomrule
\end{longtable}
\end{center}

\subsection*{Frontend API Routes (SvelteKit)}
\begin{center}
\begin{longtable}{@{}L{1.5cm}L{4.5cm}L{6cm}@{}}
\toprule
\textbf{Method} & \textbf{Endpoint} & \textbf{Description} \\
\midrule
\endhead
POST & \code{/api/casar-disciplinas} & Proxies discipline matching \\
GET  & \code{/health}                & Frontend health check \\
GET  & \code{/sitemap.xml}           & SEO sitemap \\
\bottomrule
\end{longtable}
\end{center}

% =====================================================================
\section{Authentication \& Security}
% =====================================================================

\subsection{Authentication Methods}
\begin{enumerate}[noitemsep]
    \item \textbf{Email/Password} — Standard Supabase Auth signup and login
    \item \textbf{Google OAuth} — Single sign-on via Google
    \item \textbf{Anonymous Mode} — Browse flowcharts without an account (limited features)
\end{enumerate}

\subsection{Auth Flow}
\begin{lstlisting}
Client (LoginForm) -> Supabase Auth (signIn/signUp/signInWithOAuth)
-> Receive session token -> authService.databaseSearchUser()
-> Query users table -> Set authStore -> Navigate to protected route
\end{lstlisting}

\subsection{Route Protection}
\begin{itemize}[noitemsep]
    \item \textbf{Server-side:} \code{hooks.server.ts} validates the Supabase session
    on every request to protected routes. Invalid sessions redirect to
    \code{/login?redirect=<path>}.
    \item \textbf{Client-side:} \code{authGuard.ts} checks auth state before navigation.
\end{itemize}

\subsection{Backend Authorization}

Protected backend endpoints (\code{upload-dados-fluxograma}, \code{delete-fluxograma})
require:
\begin{enumerate}[noitemsep]
    \item \code{Authorization} header with Bearer token (Supabase session)
    \item \code{User-ID} header
    \item Token validation against Supabase auth
    \item Email match verification between token and database
\end{enumerate}

\subsection{Row Level Security}

All Supabase tables use RLS policies. Academic data is read-only for authenticated
users. User-specific data (users, dados\_users) is restricted to the owning user.

% =====================================================================
\section{AI Assistant Integration}
% =====================================================================

\subsection{RAGFlow}

The AI assistant uses \textbf{RAGFlow}, an external AI agent service, to provide
discipline recommendations.

\noindent\textbf{Configuration:}
\begin{itemize}[noitemsep]
    \item \code{RAGFLOW\_API\_KEY} — API authentication
    \item \code{RAGFLOW\_BASE\_URL} — RAGFlow server URL
    \item \code{RAGFLOW\_AGENT\_ID} — Specific agent ID
\end{itemize}

\subsection{Flow}
\begin{lstlisting}
User submits subject name -> Backend normalizes text (remove accents, uppercase)
-> Create RAGFlow session -> Call RAGFlow completions API
-> Parse AI response -> Format as Markdown ranking
-> Return to client
\end{lstlisting}

\subsection{Response Format}

The AI returns a ranked list of recommended disciplines with:
\begin{itemize}[noitemsep]
    \item Discipline name and code
    \item Responsible department
    \item Relevance score (0--100)
    \item Justification text
    \item Syllabus excerpt
\end{itemize}

\subsection{Data Pipeline for AI}

Course data is formatted and uploaded to RAGFlow via
\code{coleta\_dados/scraping/formatar\_para\_ragflow.py}, which structures syllabus
and curriculum data for the AI knowledge base.

\end{document}
